\documentclass[11pt,a4paper]{article}

\usepackage[margin=1in]{geometry}
\usepackage{amsmath,amssymb,mathtools}

\newcommand{\Names}{\mathit{N}}
\newcommand{\Actor}{\mathit{Actor}}
\newcommand{\Role}{\mathit{Role}}
\newcommand{\Agent}{\mathit{Agent}}
\newcommand{\IE}{\mathit{IE}}
\newcommand{\IEA}{\mathit{IE}_{A}}
\newcommand{\IED}{\mathit{IE}_{D}}
\newcommand{\GoT}{\mathit{GoT}}
\newcommand{\RefLink}{\mathit{RefLink}}
\newcommand{\ConLink}{\mathit{ConLink}}
\newcommand{\DepenLink}{\mathit{DepenLink}}
\newcommand{\RefType}{\mathit{RefType}}
\newcommand{\ContribValue}{\mathit{ContribValue}}
\newcommand{\Par}{\mathit{Par}}
\newcommand{\Ch}{\mathit{Ch}}
\newcommand{\Leaf}{\mathit{Leaf}}
\newcommand{\GoalType}{\mathit{GoalType}}
\newcommand{\gtype}{\mathit{gtype}}
\newcommand{\BoolExpr}{\mathit{BoolExpr}}
\newcommand{\activation}{\mathit{activation}}
\newcommand{\condition}{\mathit{condition}}
\newcommand{\pre}{\mathit{pre}}
\newcommand{\post}{\mathit{post}}

\title{Formal Definition of the Goal Model}
\author{}
\date{}

\begin{document}
\maketitle

\section{Names and Actors}

Let
\[
\Names \subseteq \mathbb{A}^{+}
\]
be the set of valid names.

The finite set of actors is
\[
\Actor \subseteq \Names,
\]
with distinguished subsets
\[
\Role \subseteq \Actor,
\qquad
\Agent \subseteq \Actor.
\]

\[
\Role  \cap \Agent = \varnothing.
\]

\section{Intentional Elements}

The set of intentional elements is
\[
IE \subseteq N.
\]

For every actor \(a\in\Actor\), let
\[
Want_a
\subseteq
\{\,a\to ie \mid ie\in IE\,\}
\]
be the finite set of intentional elements declared in \(a\).

The set of actor-owned intentional elements is
\[
IE_A
=
\{
ie\in IE
\mid
\exists a\in\Actor:\ a\to ie\in Want_a
\}.
\]

Let \(IE_D\) be the finite set of dependency-owned intentional elements.

The complete set of intentional elements is
\[
IE
=
IE_A \uplus IE_D.
\]

The intentional elements are partitioned into goals, tasks, and qualities:
\[
IE
=
G \uplus T \uplus Q.
\]

For convenience, define
\[
GoT
=
G \uplus T.
\]
\section{Goal Types}

The supported goal types are
\[
\GoalType
=
\{
\mathit{Achieve},
\mathit{Maintain},
\mathit{Sustain}
\}.
\]

Each goal has exactly one goal type:
\[
\gtype
:
G
\rightarrow
\GoalType.
\]


\section{Links}

The finite set of links is
\[
L
=
\RefLink
\cup
\ConLink
\cup
\DepenLink.
\]

Refinement links connect goal-or-task elements belonging to the same actor:
\[
\RefLink
\subseteq
\GoT\times\GoT.
\]

Contribution links connect a goal-or-task element to a quality of the same actor:
\[
\ConLink
\subseteq
\GoT\times Q.
\]

Dependency links connect intentional elements owned by different actors:
\[
\DepenLink
\subseteq
\left\{
(e_1,e_2)\in\IEA\times\IEA
\;\middle|
(a \to e_1 )\in Want_a
\land
(b \to e_2)\in Want_b
\right\}.
\]

\section{AND/OR Refinement}

For every refinement link
\[
(c,p)\in\RefLink,
\]
\(c\) is the child and \(p\) is the parent.

The direct top-down refinement relation is defined by
\[
p \rightsquigarrow c
\quad\Longleftrightarrow\quad
(c,p)\in\RefLink.
\]

For every \(e\in\GoT\), its parents and children are
\[
\Par(e)
=
\{
p\in\GoT
\mid
(e,p)\in\RefLink
\},
\]
\[
\Ch(e)
=
\{
c\in\GoT
\mid
(c,e)\in\RefLink
\}.
\]

An element is a leaf iff it has no children:
\[
\Leaf(e)
\quad\Longleftrightarrow\quad
\Ch(e)=\varnothing.
\]

The refinement types are
\[
\RefType
=
\{
\mathit{AND},
\mathit{OR}
\}.
\]

The refinement-type function is
\[
R
:
\RefLink
\rightarrow
\RefType.
\]

Children of the same parent use the same refinement type:
\[
\forall e,e_1,e_2\in\GoT:
\Bigl(
(e_1,e)\in\RefLink
\land
(e_2,e)\in\RefLink
\Bigr)
\Rightarrow
R(e_1,e)=R(e_2,e).
\tag{WF-REF}
\]



\section{Contribution Links}

The contribution values are
\[
\ContribValue
=
\{
\mathit{Make},
\mathit{Help},
\mathit{Hurt},
\mathit{Break}
\}.
\]

The contribution function is
\[
C
:
\ConLink
\rightarrow
\ContribValue.
\]

Hence, for every
\[
(e,q)\in\ConLink,
\]
\[
C(e,q)
\in
\{
\mathit{Make},
\mathit{Help},
\mathit{Hurt},
\mathit{Break}
\}.
\]

Contribution branches with different contribution values do not share refinement descendants:

\[
\forall q\in Q,\ e_1,e_2\in\GoT:
\Bigl(
(e_1,q)\in\ConLink
\land
(e_2,q)\in\ConLink
\land
C(e_1,q)\neq C(e_2,q)
\Bigr)
\Rightarrow
\]
\[
\left(
\{c\mid e_1\rightsquigarrow^{*}c\}
\cap
\{c\mid e_2\rightsquigarrow^{*}c\}
=
\varnothing
\right).\tag{WF-CON}
\]

\section{Dependency Links}

Each dependency link owns exactly one intentional element that does not belong to an actor:
\[
E
:
\DepenLink
\rightarrow
\IED.
\]

Thus, for every
\[
(e_1,e_2)\in\DepenLink,
\]
\[
E(e_1,e_2)\in\IED
\]
\[
D(e_1,e_2)
\in
\{
\mathit{AND},
\mathit{OR},
\}.
\]


\section{Is-a and Participant-in Links}

Actor association links are binary relations between actors:
\[
ActorLink \subseteq Actor \times Actor.
\]

For
\[
(a,b)\in ActorLink,
\]
For an  link,
\[
(a,b)\in ActorLink
\]
means that \(a\) is a specialization of \(b\); we write
\[
a \prec_A b.
\]

Let
\[
\preceq_A
\]
denote the reflexive-transitive closure of \(\prec_A\).

\[
Want_a^{*}
=
\bigcup_{\{b\in Actor\mid a\preceq_A b\}}
Want_b.
\]


\section{Goal-Model State Expansion}

Let \(\sigma_M\) be an ACL state. 
The corresponding structural state of the goal model is determined as follows.

\medskip
\noindent
\textbf{(i) Actor instances.}

For every \(a\in\Role\),
\[
\sigma_a
=
\sigma_{\mathrm{Class}}(a).
\]

\medskip
\noindent
\textbf{(ii) Intentional-element instances.}

with \[
(e \to x)\in Want_a = e^{x},
\]  

For every \(a\in\Role\),
\[
\sigma_{Want_a}
=
\left\{
x\to e^{x}
\;\middle|\;
x\in\sigma_a
\land
\exists b\in\Actor:
(b\to e)\in Want_a^{*}
\right\}.
\]


\medskip
\noindent
\textbf{(iii) link instances.}

For every
\[
l = (c,p)\in \RefLink
\cup
\ConLink,
\]

the corresponding link instances are
\[
\sigma_{\L}(l)
=
\left\{
(c^{x},p^{x})
\;\middle|\;
x\in\sigma_a
\right\}.
\]

\[
\sigma_{\L}
=
\bigcup_{l\in \RefLink
\cup
\ConLink}
\sigma_{\L}(l),
\]


\medskip
\noindent
\textbf{(iv) Dependency instances.}

Let
\[
d=(e_1,e_2)\in\DepenLink,
\]

induces one instance of the dependency-owned intentional element
\(E(d)\in\IED\):
\[
\sigma_{E(d)}
=
\left\{
E(d)^{x,y}
\;\middle|\;
x\in\sigma_a
\land
y\in\sigma_b
\right\}.
\]

\[
\sigma_{\IED}
=
\bigcup_{d\in\DepenLink}
\sigma_{E(d)},
\]

and

\[
\sigma_{\DepenLink}(d)
=
\left\{
(e1^{x}, E(d)^{x,y},e2^{x})
\;\middle|\;
x\in\sigma_a
\land
y\in\sigma_b
\right\}.
\]

\[
\sigma_{\L}
=
\bigcup_{l\in \DepenLink}
\sigma_{\L}(l),
\]


\section{Goal-Model Marking and Evaluation}

Let
\[
\Delta
=
\{
(?,?),
(\top,\bot),
(\top,\top)
\}
\]
be the set of marking values for goals and tasks.

The value
\[
(?,?)
\]
denotes an element that has not yet been activated or determined,

\[
(\top,\bot)
\]
denotes an activated but not fulfilled element, and

\[
(\top,\top)
\]
denotes a fulfilled element.


\subsection{State-Based OCL Evaluation}

Let
\[
\sigma_M^i
\]
be the ACL state at step \(i\).

For every goal instance \(g^x\), define
\[
eval_{\sigma_M^i}(g^x)
\in
\{\top,\bot\},
\]
where
\[
eval_{\sigma_M^i}(g^x)=\top
\]
iff the OCL condition of \(g\), evaluated in
\(\sigma_M^i\) with actor instance \(x\), is true.

For every task instance \(t^x\), define
\[
evalPre_{\sigma_M^i}(t^x)
\in
\{\top,\bot\},
\]
where
\[
evalPre_{\sigma_M^i}(t^x)=\top
\]
iff the precondition of \(t\), evaluated in
\(\sigma_M^i\) with actor instance \(x\), is true.

Similarly,
\[
evalPost_{\sigma_M^i}(t^x)
\in
\{\top,\bot\},
\]
where
\[
evalPost_{\sigma_M^i}(t^x)=\top
\]
iff the postcondition of \(t\), evaluated in
\(\sigma_M^i\) with actor instance \(x\), is true.


\subsection{Goal Evaluation}

For every goal instance \(g^x\), let
\[
evalGoal_i(g^x)\in\Delta
\]
denote its marking at state \(\sigma_M^i\).

\medskip
\noindent
\textbf{(i) Maintain goal.}

If
\[
\gtype(g)=\mathit{Maintain},
\]
then
\[
evalGoal_i(g^x)
=
\begin{cases}
(\top,\top),
&
eval_{\sigma_M^i}(g^x)=\top,
\\[1mm]
(\top,\bot),
&
eval_{\sigma_M^i}(g^x)=\bot.
\end{cases}
\]

\medskip
\noindent
\textbf{(ii) Achieve goal.}

If
\[
\gtype(g)=\mathit{Achieve},
\]
then
\[
evalGoal_i(g^x)
=
\begin{cases}
(\top,\top),
&
evalGoal_{i-1}(g^x)=(\top,\top),
\\[1mm]
(\top,\top),
&
eval_{\sigma_M^i}(g^x)=\top,
\\[1mm]
(?,?),
&
\text{otherwise}.
\end{cases}
\]

Hence, once an Achieve goal is fulfilled, it remains fulfilled.

\medskip
\noindent
\textbf{(iii) Sustain goal.}

If
\[
\gtype(g)=\mathit{Sustain},
\]
then
\[
evalGoal_i(g^x)
=
\begin{cases}
(\top,\bot),
&
evalGoal_{i-1}(g^x)=(\top,\bot),
\\[1mm]
(\top,\bot),
&
evalGoal_{i-1}(g^x)=(\top,\top)
\land
eval_{\sigma_M^i}(g^x)=\bot,
\\[1mm]
(\top,\top),
&
eval_{\sigma_M^i}(g^x)=\top,
\\[1mm]
(?,?),
&
\text{otherwise}.
\end{cases}
\]

Thus, a Sustain goal remains unknown until it is first fulfilled.
After it has been fulfilled, a later violation changes its marking to
\((\top,\bot)\), which is preserved thereafter.


\subsection{Task Evaluation}

For every task instance \(t^x\), let
\[
evalTask_i(t^x)\in\Delta
\]
denote its marking at state \(\sigma_M^i\).

The task marking is defined by
\[
evalTask_i(t^x)
=
\begin{cases}
(\top,\top),
&
evalTask_{i-1}(t^x)=(\top,\top),
\\[1mm]
(\top,\top),
&
evalTask_{i-1}(t^x)=(\top,\bot)
\land
evalPost_{\sigma_M^i}(t^x)=\top,
\\[1mm]
(\top,\bot),
&
evalTask_{i-1}(t^x)=(\top,\bot),
\\[1mm]
(\top,\bot),
&
evalTask_{i-1}(t^x)=(?,?)
\land
evalPre_{\sigma_M^i}(t^x)=\top,
\\[1mm]
(?,?),
&
\text{otherwise}.
\end{cases}
\]

Hence, the task evolves according to
\[
(?,?)
\xrightarrow{\;pre\;}
(\top,\bot)
\xrightarrow{\;post\;}
(\top,\top).
\]

A postcondition that holds before the corresponding precondition has
been activated does not change the task marking.


\subsection{Initial Evaluation}

Before the first ACL state is evaluated, every goal and task instance
has the initial value
\[
evalGoal_{-1}(g^x)=(?,?),
\]
and
\[
evalTask_{-1}(t^x)=(?,?).
\]

\subsection{ACL-State-Path Semantics}

The goal model does not introduce an independent state space.  For a finite
ACL state path
\[
\bar{\sigma}_M
=
\langle\sigma_M^0,\ldots,\sigma_M^n\rangle,
\]
the iStar structural state and marking at position $i$ are derived from
$\sigma_M^i$ by the expansion and evaluation rules above.  Thus
\[
\sigma_I^i
=
Expand_I(GM_M,\sigma_M^i),
\]
and $\sigma_I^i$ is not chosen independently of $\sigma_M^i$.

Only leaf goals may own a directly evaluated condition:
\[
\forall g\in G:
\quad
g\in\operatorname{dom}(\condition)
\Longrightarrow
\Ch(g)=\varnothing.
\tag{WF-GOAL-CONDITION}
\]
A non-leaf goal is evaluated exclusively by its AND/OR refinement.  This
prevents a parent condition from bypassing or contradicting its children.


\subsection{Refinement Evaluation}

Only leaf goals and tasks may be directly specified by OCL conditions.
A non-leaf element is evaluated from its refinement children.

Let
\[
\Ch_{\sigma}(p^x)
=
\{
e_1^x,\ldots,e_n^x
\}
\]
be the instantiated children of \(p^x\).

Let
\[
eval_i(e^x)
=
\begin{cases}
evalGoal_i(e^x),
&
e\in G,
\\
evalTask_i(e^x),
&
e\in T.
\end{cases}
\]

\medskip
\noindent
\textbf{(i) AND refinement.}

If the refinement of \(p\) is AND, then
\[
evalAND_i(p^x)
=
\begin{cases}
(\top,\top),
&
\forall e^x\in\Ch_{\sigma}(p^x):
eval_i(e^x)=(\top,\top),
\\[1mm]
(\top,\bot),
&
\exists e^x\in\Ch_{\sigma}(p^x):
eval_i(e^x)=(\top,\bot),
\\[1mm]
(?,?),
&
\text{otherwise}.
\end{cases}
\]

\medskip
\noindent
\textbf{(ii) OR refinement.}

If the refinement of \(p\) is OR, then
\[
evalOR_i(p^x)
=
\begin{cases}
(\top,\top),
&
\exists e^x\in\Ch_{\sigma}(p^x):
eval_i(e^x)=(\top,\top),
\\[1mm]
(\top,\bot),
&
\forall e^x\in\Ch_{\sigma}(p^x):
eval_i(e^x)=(\top,\bot),
\\[1mm]
(?,?),
&
\text{otherwise}.
\end{cases}
\]

\section{Bounded ACL--BPMN--iStar Consistency}

Let the root goals be
\[
G_{root}
=
\{g\in G\mid\Par(g)=\varnothing\}.
\]
Root tasks are not verdict targets.  Non-root goals and tasks contribute only
through refinement propagation to a member of $G_{root}$.

For an actor instance $x$ present in the final ACL state, define
\[
Achieves_I(\bar{\sigma}_M)
\Longleftrightarrow
\forall g\in G_{root}:
Inst_{owner(g)}(\sigma_M^n)\neq\varnothing
\land
\forall x\in Inst_{owner(g)}(\sigma_M^n):
evalGoal_n(g^x)=(\top,\top).
\]

Let $Runs_{k,s}(B)$ be the maximal BPMN executions admitted by loop bound $k$
and snapshot bound $s$, and let
\[
Path_M(\rho,\beta)
\]
be the ACL state paths inside boundary $\beta$ that satisfy every BPMN
transition contract along $\rho$.  The realizable bounded runs are
\[
Runs^M_{k,s,\beta}(B)
=
\{\rho\in Runs_{k,s}(B)\mid Path_M(\rho,\beta)\neq\varnothing\}.
\]
Let $Admissible_{k,s}(\rho)$ mean that $\rho$ reaches an End configuration or
a boundary cutoff, rather than a structural deadlock.
Full bounded consistency is
\[
Consistent_{k,s,\beta}(M,B,GM_M)
\Longleftrightarrow
Runs^M_{k,s,\beta}(B)\neq\varnothing
\land
\forall\rho\in Runs^M_{k,s,\beta}(B):
Admissible_{k,s}(\rho)
\land
\exists\bar{\sigma}_M\in Path_M(\rho,\beta):
Achieves_I(\bar{\sigma}_M).
\]

Weak bounded consistency is
\[
WeakConsistent_{k,s,\beta}(M,B,GM_M)
\Longleftrightarrow
\exists\rho\in Runs^M_{k,s,\beta}(B),\
\exists\bar{\sigma}_M\in Path_M(\rho,\beta):
Admissible_{k,s}(\rho)
\land
Achieves_I(\bar{\sigma}_M).
\]
It is reported only when full consistency is false.  Inconsistency is the
negation of weak consistency.

An explanatory activity mapping may be returned as a partial function
\[
Map_{BI}:Activity\rightharpoonup
\{e\in GoT\mid\Leaf(e)\}.
\]
This mapping records contract/name correspondence for inspection; the
consistency verdict is determined by the ACL-state-path formulas above and
does not assume that the inferred mapping is an additional semantic axiom.


\end{document}
